%-------------------------------------------------------------------------------
%   ӨГӨГДЛИЙН ГАРАЛ ҮҮСЛИЙН СИСТЕМ - Бүрэн Төслийн Тайлан
%-------------------------------------------------------------------------------

\section{Өгөгдлийн Гарал Үүслийн Систем}
\label{sec:data_lineage_system}

%-------------------------------------------------------------------------------
% БҮЛГИЙН АГУУЛГА
%-------------------------------------------------------------------------------
\subsection*{Агуулга}
\begin{enumerate}
    \item Уран зохиолын тойм
    \item Удиртгал
    \item Өгөгдлийн гарал үүсэл
    \item Хамрах хүрээ
    \item Өгөгдлийн гарал үүслийн чухал тодорхойлолтууд
    \item Зохиомж
    \item Төслийн суурь мэдээлэл
    \item Санал болгож буй шийдэл
    \item Шаардлагын жагсаалт
    \item UML диаграммууд
    \item Хэрэгжүүлэлт
    \item Дүгнэлт
    \item Ном зүй
\end{enumerate}

\hrulefill

%-------------------------------------------------------------------------------
% 1. УРАН ЗОХИОЛЫН ТОЙМ
%-------------------------------------------------------------------------------
\subsection{Уран зохиолын тойм}
\label{subsec:literature_survey}

Өгөгдлийн гарал үүслийн салбар нь сүүлийн арван жилд зохицуулалтын шаардлага болон байгууллагын орчинд өгөгдлийн засаглалын хэрэгцээнээс үүдэн ихээхэн хөгжсөн байна. Энэ хэсэгт орчин үеийн өгөгдлийн гарал үүслийн системүүдийг бүрдүүлсэн гол эрдэм шинжилгээний болон салбарын бүтээлүүдийг авч үзнэ.

\textbf{Thompson ба Garcia (2019)} метадата удирдлагын фреймворкуудын талаар иж бүрэн судалгаа хийж, автоматжуулсан гарал үүслийн хяналт нь гарын авлагын баримтжуулалтын ажлыг 60\% хүртэл бууруулдаг болохыг онцолсон. Тэдний судалгаагаар системчилсэн гарал үүслийн хяналтыг хэрэгжүүлсэн байгууллагууд өгөгдлийн чанарын асуудлыг 40\%-иар бууруулсан байна.

\textbf{Chen нар (2020)} өгөгдлийн гарал үүслийн дүрслэлд граф дээр суурилсан арга барилыг санал болгож, уламжлалт хүснэгт хэлбэрийн дүрслэл нь орчин үеийн өгөгдлийн экосистемийн нарийн төвөгтэй байдлыг илэрхийлж чаддаггүй гэж маргасан. Тэдний фреймворк нь эцсийн хэрэглэгчдийн гарал үүслийн ойлголтыг сайжруулсан шаталсан кластерийн алгоритмуудыг нэвтрүүлсэн.

\textbf{Williams ба Brown (2018)} ETL хэрэгслүүд болон гарал үүслийн системүүдийн хоорондох интеграцийн сорилтуудыг судалсан. Тэдний судалгаагаар байгууллагуудын 73\% нь олон платформ дээр тархсан гарал үүслийн мэдээлэлтэй тэмцэж байгаа нь нэгдсэн гарал үүслийн шийдлийн хэрэгцээг онцолсон.

\textbf{Kumar ба Patel (2021)} автомат SQL гарал үүсэл гаргаж авахад машин сургалтад суурилсан задлагч хөгжүүлсэн. Тэдний систем нарийн төвөгтэй query мэдэгдлүүдээс эх-зорилтот харилцааг тодорхойлоход 94\% нарийвчлалтай байсан.

\textbf{Anderson (2020)} өгөгдлийн гарал үүслийн зохицуулалтын үр нөлөөг, ялангуяа GDPR болон CCPA-ийн нийцлийн талаар судалсан. Судалгаагаар зохих гарал үүслийн баримтжуулалт нь нийцлийн аудитын хугацааг дунджаар 35\%-иар бууруулдаг гэж дүгнэсэн.

\textbf{Liu ба Yamamoto (2022)} урсгал өгөгдлийн архитектурт бодит цагийн гарал үүслийн хяналтыг судалж, нарийвчлалыг хадгалахын зэрэгцээ хоцрогдлыг багасгах үйл явдалд суурилсан гарал үүслийн барих механизмуудыг санал болгосон.

Эдгээр судалгаанууд нь санал болгож буй Өгөгдлийн Гарал Үүслийн Системийн онолын үндэс суурийг бүрдүүлж, архитектурын шийдвэр болон функцийн эрэмбэлэлтэд мэдээлэл өгдөг.

%-------------------------------------------------------------------------------
% 2. УДИРТГАЛ
%-------------------------------------------------------------------------------
\subsection{Удиртгал}
\label{subsec:dl_introduction}

Өнөөгийн өгөгдөлд суурилсан байгууллагын орчинд байгууллагууд олон өгөгдлийн сан, өгөгдлийн агуулах, ETL дамжуулах хоолой, аналитик хэрэгслүүдээс бүрдсэн улам бүр нарийн төвөгтэй өгөгдлийн экосистемийг удирддаг. Өгөгдөл эдгээр системүүдээр хэрхэн урсдагийг---эх үүсвэрээс хэрэглээ хүртэл---ойлгох нь өгөгдлийн засаглал, зохицуулалтын нийцэл, үйл ажиллагааны үр ашгийн чухал шаардлага болсон.

Өгөгдлийн гарал үүсэл гэдэг нь өгөгдлийн амьдралын мөчлөг буюу түүний гарал үүсэл, хувиргалт, эцсийн очих газрыг хэлнэ. Найдвартай өгөгдлийн гарал үүслийн систем нь өгөгдлийн хөдөлгөөнд ил тод байдлыг хангаж, оролцогч талуудад дараах боломжийг олгоно:

\begin{itemize}
    \item \textbf{Өгөгдлийн чанарыг хангах:} Өгөгдлийн асуудлуудыг түргэн засварлахын тулд эх үүсвэр хүртэл нь мөшгих.
    \item \textbf{Зохицуулалтын нийцлийг хадгалах:} Аудитор болон зохицуулалтын байгууллагуудад зориулж өгөгдлийн урсгалыг баримтжуулах.
    \item \textbf{Нөлөөллийн шинжилгээг дэмжих:} Өөрчлөлт хийхээс өмнө доод түвшний нөлөөллийг ойлгох.
    \item \textbf{Өгөгдөл олохыг идэвхжүүлэх:} Аналистуудад холбогдох өгөгдлийн хөрөнгийг үр ашигтай олоход туслах.
\end{itemize}

Энэ төсөл нь олон төрлийн өгөгдлийн орчинд өгөгдлийн гарал үүслийн мэдээллийг автоматаар барьж авах, дүрслэх, удирдахад зориулагдсан иж бүрэн веб-д суурилсан Өгөгдлийн Гарал Үүслийн Хэрэгслийг танилцуулж байна. Систем нь SQL скриптүүд, өгөгдлийн сангийн схемүүд, ETL тохиргоонуудыг задлан шинжилж нэгдсэн гарал үүслийн граф бүтээх замаар автоматжуулсан гарал үүслийн хяналтын өсөн нэмэгдэж буй хэрэгцээг хангана.

Санал болгож буй шийдэл нь байгууллагынхаа өгөгдлийн орчинд тодорхой харагдах байдал шаарддаг өгөгдлийн инженерүүд, өгөгдлийн аналистууд, засаглалын багуудад зориулагдсан. Зөн совинтой дүрслэл болон хүчирхэг хайлтын чадамжуудыг хангаснаар систем нь нарийн төвөгтэй метадата харилцааг үйл ажиллагаанд хэрэглэгдэх мэдээлэл болгон хувиргана.

Энэ тайлан нь системийн зохиомж, хэрэгжүүлэлт, байршуулалтын үзэл баримтлалыг дэлгэрэнгүй тайлбарлаж, Өгөгдлийн Гарал Үүслийн Системийн бүрэн техникийн төлөвлөгөөг өгнө.

%-------------------------------------------------------------------------------
% 3. ӨГӨГДЛИЙН ГАРАЛ ҮҮСЭЛ
%-------------------------------------------------------------------------------
\subsection{Өгөгдлийн гарал үүсэл}
\label{subsec:data_lineage_concept}

Өгөгдлийн гарал үүсэл гэдэг нь өгөгдлийг эх үүсвэрээс янз бүрийн хувиргалтуудаар дамжуулан эцсийн очих газар хүртэл мөшгих үйл явц юм. Энэ нь байгууллагын мэдээллийн архитектур дотор өгөгдлийн хөдөлгөөний иж бүрэн зураглалыг бий болгож, үндсэн асуултуудад хариулт өгдөг: \textit{Энэ өгөгдөл хаанаас ирсэн бэ? Хэрхэн хувиргагдсан бэ? Хаана ашиглагддаг вэ?}

\subsubsection{Өгөгдлийн гарал үүслийн төрлүүд}

\paragraph{Урагшлах гарал үүсэл (Нөлөөллийн шинжилгээ)}
Урагшлах гарал үүсэл нь өгөгдлийг эх үүсвэрээс очих газар хүртэл мөшгиж, бүх доод түвшний хамаарлуудыг илрүүлдэг. Энэ нь өөрчлөлтийн нөлөөллийн шинжилгээнд маш чухал---эх хүснэгтийг өөрчлөхөд аль тайлан, хянах самбар эсвэл програмуудад нөлөөлөхийг ойлгох.

\paragraph{Ухрах гарал үүсэл (Үндсэн шалтгааны шинжилгээ)}
Ухрах гарал үүсэл нь өгөгдлийг одоогийн байршлаас эх үүсвэр хүртэл нь буцааж мөшгидөг. Энэ нь өгөгдлийн чанарын асуудал илэрсэн үед үндсэн шалтгааны шинжилгээг боломжтой болгож, инженерүүдэд асуудлууд хаанаас үүссэнийг яг таг тодорхойлох боломжийг олгодог.

\paragraph{Хэвтээ гарал үүсэл}
Хэвтээ гарал үүсэл нь ижил архитектурын түвшинд өгөгдлийн хөдөлгөөнийг харуулдаг, жишээлбэл өгөгдлийн агуулах доторх хүснэгтээс хүснэгт рүү хувиргалтууд.

\subsubsection{Гарал үүслийн нарийвчлал}

Өгөгдлийн гарал үүслийг олон нарийвчлалын түвшинд барьж авч болно:

\begin{itemize}
    \item \textbf{Системийн түвшин:} Програмууд болон өгөгдлийн сангуудын хоорондын холболтууд.
    \item \textbf{Өгөгдлийн олонлогийн түвшин:} Хүснэгтүүд, файлууд, тайлангуудын хоорондын харилцаа.
    \item \textbf{Баганын түвшин:} Тус тусдаа талбарын хувиргалтуудын нарийвчилсан хяналт.
\end{itemize}

Орчин үеийн өгөгдлийн гарал үүслийн системүүд баганын түвшний гарал үүсэлд улам бүр анхаарч байгаа нь зохицуулалтын нийцэл болон дэлгэрэнгүй нөлөөллийн шинжилгээнд шаардлагатай нарийвчлалыг хангадаг. Бидний санал болгож буй систем нь гурван нарийвчлалын түвшинг бүгдийг дэмждэг бөгөөд хамгийн их ашиг тустай байхын тулд баганын түвшний хяналтад онцгой анхаарал хандуулдаг.

%-------------------------------------------------------------------------------
% 4. ХАМРАХ ХҮРЭЭ
%-------------------------------------------------------------------------------
\subsection{Хамрах хүрээ}
\label{subsec:dl_scope}

Өгөгдлийн Гарал Үүслийн Системийн төсөл нь автоматжуулсан өгөгдлийн гарал үүслийн хяналтад зориулсан веб-д суурилсан програмын зохиомж, хөгжүүлэлт, байршуулалтыг хамарна. Хамрах хүрээ нь төслийн гаралт болон функцийн тодорхой хязгаарыг тодорхойлно.

\subsubsection{Хамрах хүрээнд багтах}

\begin{itemize}
    \item \textbf{Метадата гаргаж авалт:} Дэмжигдсэн өгөгдлийн сангийн платформуудаас (PostgreSQL, MySQL, Oracle, SQL Server) SQL DDL/DML мэдэгдлүүд, хадгалагдсан процедурууд, view тодорхойлолтуудыг автоматаар задлан шинжлэх.
    
    \item \textbf{ETL интеграци:} Apache Airflow DAG-ууд болон тусгай Python скриптүүд зэрэг түгээмэл хэрэгслүүдээс ETL ажлын тодорхойлолтуудыг задлан шинжлэх.
    
    \item \textbf{Гарал үүслийн дүрслэл:} Томруулах, гүйлгэх, шүүх чадвартай өгөгдлийн урсгалын интерактив граф-д суурилсан дүрслэл.
    
    \item \textbf{Хайлт ба олох:} Талын шүүлтүүрийн сонголттой бүх метадата нэгжүүдээр бүрэн текст хайлт.
    
    \item \textbf{Хэрэглэгчийн удирдлага:} Үзэгч, засварлагч, администраторуудыг дэмждэг үүрэгт суурилсан хандалтын хяналт.
    
    \item \textbf{Экспортын функц:} Олон форматаар (PDF, PNG, JSON, CSV) гарал үүслийг экспортлох.
    
    \item \textbf{API хандалт:} Программчлалын гарал үүслийн асуулга болон интеграцид зориулсан RESTful API.
\end{itemize}

\subsubsection{Хамрах хүрээнд багтахгүй}

\begin{itemize}
    \item Бодит цагийн урсгал гарал үүсэл барих (ирээдүйн хувилбарт төлөвлөгдсөн).
    \item Өгөгдлийн чанарын хяналт болон сэрэмжлүүлэг.
    \item Автоматжуулсан өгөгдлийн ангилал болон мэдрэмтгий байдлын шошголол.
    \item Арилжааны BI хэрэгслүүдтэй (Tableau, Power BI) интеграци.
    \item Гар утасны програм хөгжүүлэлт.
\end{itemize}

Энэ хамрах хүрээ нь ирээдүйн сайжруулалтын суурийг тавихын зэрэгцээ төвлөрсөн хөгжүүлэлтийг хангана.

%-------------------------------------------------------------------------------
% 5. ӨГӨГДЛИЙН ГАРАЛ ҮҮСЛИЙН ЧУХАЛ ТОДОРХОЙЛОЛТУУД
%-------------------------------------------------------------------------------
\subsection{Өгөгдлийн гарал үүслийн чухал тодорхойлолтууд}
\label{subsec:dl_definitions}

Өгөгдлийн гарал үүслийг ойлгоход гол нэр томъёотой танилцах шаардлагатай. Энэ хэсэгт энэ тайланд ашигласан чухал ойлголтуудыг тодорхойлно.

\begin{description}
    \item[Өгөгдлийн хөрөнгө (Data Asset)] Хүснэгт, view, файл, тайлан, ETL ажлууд зэрэг өгөгдөл агуулсан эсвэл боловсруулдаг аливаа нэгж.
    
    \item[Эх үүсвэр (Source)] Гарал үүслийн харилцаан дахь өгөгдлийн эх цэг. Эх үүсвэр нь нэг буюу түүнээс олон зорилт руу урсдаг өгөгдлийг хангана.
    
    \item[Зорилт (Target)] Гарал үүслийн харилцаан дахь өгөгдлийн очих газар. Зорилт нь нэг буюу түүнээс олон эх үүсвэрээс өгөгдөл хүлээн авна.
    
    \item[Хувиргалт (Transformation)] Шүүлтүүр, нэгтгэл, холболт, тооцоолол зэрэг өгөгдлийг эх үүсвэрээс зорилт руу шилжих үед өөрчилдөг аливаа үйлдэл.
    
    \item[Зангилаа (Node)] Өгөгдлийн хөрөнгийг төлөөлөх гарал үүслийн графын орой. Зангилаа бүр хөрөнгийн шинж чанарыг тодорхойлсон метадата агуулна.
    
    \item[Ирмэг (Edge)] Хоёр зангилааны хооронд өгөгдлийн урсгалыг төлөөлөх гарал үүслийн графын холболт. Ирмэгүүд нь чиглэлтэй бөгөөд урсгалын чиглэлийг заана.
    
    \item[Дээд урсгал (Upstream)] Өгөгдсөн хөрөнгөд өгөгдөл нийлүүлдэг бүх өгөгдлийн хөрөнгүүд. Дээд урсгалын шинжилгээ нь өгөгдлийн эх үүсвэрүүдийг илрүүлдэг.
    
    \item[Доод урсгал (Downstream)] Өгөгдсөн хөрөнгөөс өгөгдөл хүлээн авдаг бүх өгөгдлийн хөрөнгүүд. Доод урсгалын шинжилгээ нь өгөгдлийн хэрэглэгчдийг илрүүлдэг.
    
    \item[Метадата (Metadata)] Схемийн тодорхойлолт, тайлбар, эзэмшигч, техникийн шинж чанаруудыг багтаасан өгөгдлийн хөрөнгийн тухай тодорхойлолтын мэдээлэл.
    
    \item[Схем (Schema)] Баганууд, өгөгдлийн төрлүүд, хязгаарлалтууд, харилцааг тодорхойлсон өгөгдлийн сангийн объектын бүтцийн тодорхойлолт.
    
    \item[ETL (Extract, Transform, Load)] Эх үүсвэрүүдээс өгөгдөл гаргаж авах, хувиргалтуудыг хэрэгжүүлэх, үр дүнг зорилт руу ачаалах өгөгдлийн интеграцийн үйл явц.
    
    \item[DAG (Directed Acyclic Graph)] Чиглэлтэй ирмэгүүд бүхий, мөчлөггүй графын бүтэц бөгөөд өгөгдлийн дамжуулах хоолой болон гарал үүслийг төлөөлөхөд түгээмэл ашиглагддаг.
    
    \item[Баганын түвшний гарал үүсэл (Column-Level Lineage)] Тус тусдаа эх баганыг хувиргалтуудаар дамжуулан зорилтын баганатай холбодог нарийвчилсан гарал үүслийн хяналт.
\end{description}

%-------------------------------------------------------------------------------
% 6. ЗОХИОМЖ
%-------------------------------------------------------------------------------
\subsection{Зохиомж}
\label{subsec:dl_design}

Өгөгдлийн Гарал Үүслийн Систем нь өргөтгөх чадвар, засвар үйлчилгээний хялбар байдал, өргөтгөлийн боломжид зориулагдсан модульчлагдсан, микросервис-д суурилсан архитектурыг дагадаг. Энэ хэсэгт өндөр түвшний зохиомжийн зарчмууд болон архитектурын бүрэлдэхүүн хэсгүүдийг тоймлоно.

\subsubsection{Архитектурын зарчмууд}

\begin{itemize}
    \item \textbf{Үүргийн салгалт:} Бүрэлдэхүүн хэсэг бүр тодорхой хариуцлагыг хариуцаж, бие даасан хөгжүүлэлт болон тестлэлтийг боломжтой болгоно.
    
    \item \textbf{Өргөтгөх чадвар:} Плагин-д суурилсан задлагчийн архитектур нь үндсэн системийг өөрчлөхгүйгээр шинэ өгөгдлийн эх үүсвэрүүдийн дэмжлэг нэмэх боломжийг олгоно.
    
    \item \textbf{Масштаблах чадвар:} Төлөвгүй API үйлчилгээнүүд нь нэмэгдсэн метадата хэмжээг зохицуулахын тулд хэвтээ масштаблалтыг дэмждэг.
    
    \item \textbf{Аюулгүй байдлын зохиомж:} Баталгаажуулалт, зөвшөөрөл, аудит бүртгэл нь архитектурын түвшинд интеграцилагдсан.
\end{itemize}

\subsubsection{Системийн бүрэлдэхүүн хэсгүүд}

\paragraph{Танилцуулгын давхарга}
Хэрэглэгчийн интерфэйсийг хангадаг React-д суурилсан нэг хуудсан програм (SPA). Функцууд нь D3.js ашигласан интерактив гарал үүслийн графууд, метадата үзэх, удирдлагын функцуудыг багтаана.

\paragraph{API давхарга}
Бизнесийн логик, баталгаажуулалт, өгөгдлийн хандалтыг зохицуулдаг Node.js/Express-ээр бүтээгдсэн RESTful API үйлчилгээнүүд. GraphQL endpoint-ууд нь гарал үүслийн дамжуулалтад уян хатан асуулга хийх боломжийг олгоно.

\paragraph{Задлагчийн хөдөлгүүр}
SQL задлан шинжлэлд ANTLR, ETL хэрэгслүүдэд тусгай задлагчид ашигладаг Python-д суурилсан задлан шинжлэлийн фреймворк. Задлагчид нь эх үүсвэрийн бүтээгдэхүүнүүдээс гарал үүслийн харилцаа болон метадатаг гаргаж авна.

\paragraph{Граф өгөгдлийн сан}
Neo4j нь гарал үүслийн харилцааны үндсэн хадгалалт болж, үр ашигтай дамжуулалтын асуулгад уугуул графын чадамжуудыг ашигладаг. Харилцааны өгөгдлийн санд үнэтэй join шаарддаг харилцааны асуулгууд миллисекундэд биелдэг.

\paragraph{Метадата хадгалах сан}
PostgreSQL нь нэмэлт метадата, хэрэглэгчийн мэдээлэл, аудит бүртгэл, харилцааны бүтцээс ашиг хүртдэг тохиргооны өгөгдлийг хадгална.

\subsubsection{Технологийн стек}

\begin{table}[htbp]
\centering
\caption{Технологийн стекийн тойм}
\label{tab:tech_stack}
\begin{tabular}{|l|l|}
\hline
\textbf{Давхарга} & \textbf{Технологи} \\
\hline
Frontend & React, TypeScript, D3.js \\
\hline
Backend & Node.js, Express, Python \\
\hline
Граф өгөгдлийн сан & Neo4j \\
\hline
Харилцааны өгөгдлийн сан & PostgreSQL \\
\hline
Баталгаажуулалт & JWT, OAuth 2.0 \\
\hline
Контейнержуулалт & Docker, Kubernetes \\
\hline
\end{tabular}
\end{table}

%-------------------------------------------------------------------------------
% 7. ТӨСЛИЙН СУУРЬ МЭДЭЭЛЭЛ
%-------------------------------------------------------------------------------
\subsection{Төслийн суурь мэдээлэл}
\label{subsec:dl_background}

Өгөгдлийн Гарал Үүслийн Системийн төсөл нь байгууллагын өгөгдлийн удирдлагын чадамжийн чухал цоорхойгоос үүссэн. Байгууллагууд өгөгдөлд суурилсан шийдвэр гаргалтад улам бүр найддаг болохын хэрээр тэдний өгөгдлийн экосистемийн нарийн төвөгтэй байдал экспоненциалаар өссөн. Энэ хэсэгт төслийг салбарын өргөн хандлага болон байгууллагын хэрэгцээний хүрээнд тайлбарлана.

\subsubsection{Салбарын нөхцөл байдал}

Дэлхийн өгөгдлийн засаглалын зах зээл нь зохицуулалтын дарамт (GDPR, CCPA, SOX) болон өгөгдлийн хөрөнгийн стратегийн ач холбогдлоос үүдэн 2025 он гэхэд 5.7 тэрбум ам.долларт хүрэхээр төлөвлөгдөж байна. Энэ зах зээлийн дотор өгөгдлийн гарал үүслийн хэрэгслүүд нь хамгийн хурдан өсөж буй сегментүүдийн нэгийг төлөөлж, байгууллагууд гарын авлагын баримтжуулалтын практикийг орлох автоматжуулсан шийдлүүдийг хайж байна.

Гарал үүслийн баримтжуулалтын уламжлалт аргууд---хүснэгтүүд, wiki хуудсууд, хүмүүсийн мэдлэг---орчин үеийн өгөгдлийн орчинд хангалтгүй болохыг нотолсон. Гарын авлагын баримтжуулалт хурдан хуучирдаг, нийцлийн шаардлагатай нарийвчлал дутагддаг, өгөгдлийн хэмжээ нэмэгдэхэд муу масштаблагддаг.

\subsubsection{Байгууллагын сорилтууд}

Энэ төслийг санаачлахад хүргэсэн ердийн байгууллагын сорилтууд:

\begin{itemize}
    \item \textbf{Тусгаарлагдсан метадата:} Гарал үүслийн мэдээлэл олон хэрэгсэл болон багуудад тархсан, нэгдсэн харагдац байхгүй.
    
    \item \textbf{Гарын авлагын баримтжуулалтын ачаалал:} Өгөгдлийн инженерүүд цагийнхаа 20-30\%-ийг хөгжүүлэлтийн оронд баримтжуулалтад зарцуулдаг.
    
    \item \textbf{Нийцлийн эрсдэл:} Зохицуулалтын аудитын үеэр өгөгдлийн гарал үүслийг харуулж чадахгүй байдал.
    
    \item \textbf{Өөрчлөлтийн удирдлагын алдаа:} Доод түвшний нөлөөллийг ойлгохгүйгээр өөрчлөлт хийснээс үүссэн үйлдвэрлэлийн ослууд.
    
    \item \textbf{Өгөгдөл олох үр ашиггүй байдал:} Аналистууд холбогдох өгөгдлийн олонлогуудыг олж чадахгүй байгаа нь давхар ажилд хүргэдэг.
\end{itemize}

\subsubsection{Төслийн үүсэл}

Энэ төсөл нь автоматжуулалт болон дүрслэлээр дамжуулан эдгээр сорилтуудыг шийдвэрлэхэд санаачлагдсан. Шийдэл нь гарын авлагын гарал үүслийн баримтжуулалтыг устгах, өгөгдлийн урсгалд бодит цагийн харагдах байдлыг хангах, байгууллага даяарх метадатагийн эрх мэдэл бүхий эх үүсвэр болохыг зорьж байна.

%-------------------------------------------------------------------------------
% 8. САНАЛ БОЛГОЖ БУЙ ШИЙДЭЛ
%-------------------------------------------------------------------------------
\subsection{Санал болгож буй шийдэл}
\label{subsec:dl_proposed_solution}

Санал болгож буй Өгөгдлийн Гарал Үүслийн Систем нь гарал үүслийн барилтыг автоматжуулж, зөн совинтой дүрслэл хангаж, одоо байгаа өгөгдлийн дэд бүтэцтэй интеграцилагдсан иж бүрэн веб-д суурилсан платформ юм. Энэ хэсэгт шийдлийн гол чадамжууд болон ялгарах онцлогуудыг тайлбарлана.

\subsubsection{Шийдлийн тойм}

Систем нь метадата цуглуулалт, задлан шинжлэл, хадгалалт, танилцуулгын тасралтгүй мөчлөгөөр ажилладаг:

\begin{enumerate}
    \item \textbf{Цуглуулалт:} Автоматжуулсан агентууд тохируулагдсан өгөгдлийн эх үүсвэрүүдийг сканнерддаг, DDL скриптүүд, асуулгын бүртгэлүүд, ETL тодорхойлолтуудыг гаргаж авна.
    
    \item \textbf{Задлан шинжлэл:} Мэргэжлийн задлагчид цуглуулсан бүтээгдэхүүнүүдийг шинжилж, нэгжүүд, харилцаанууд, хувиргалтуудыг тодорхойлно.
    
    \item \textbf{Хадгалалт:} Задлан шинжилсэн гарал үүслийг харилцааны дамжуулалтад оновчлогдсон граф өгөгдлийн санд хадгална.
    
    \item \textbf{Танилцуулга:} Хэрэглэгчид интерактив графууд болон хайлтыг онцолсон зөн совинтой веб интерфэйсээр гарал үүсэлтэй харилцана.
\end{enumerate}

\subsubsection{Гол онцлогууд}

\paragraph{Автоматжуулсан гарал үүслийн гаргаж авалт}
Систем нь гарын авлагын оролт шаардахгүйгээр хүснэгт болон баганын түвшний гарал үүслийг гаргаж авахын тулд SQL мэдэгдлүүдийг автоматаар задлан шинжилдэг. CTE-үүд, дэд асуулгууд, цонхны функцууд зэрэг нарийн төвөгтэй асуулгуудын дэмжлэг нь иж бүрэн хамрах хүрээг хангана.

\paragraph{Интерактив дүрслэл}
Хүчний чиглэлтэй графын дүрслэл нь хэрэглэгчдэд гарал үүслийн харилцааг зөн совиноор судлах боломжийг олгоно. Функцууд нь дараах зүйлсийг багтаана:
\begin{itemize}
    \item Томруулах болон гүйлгэх навигаци
    \item Зангилааны бүлгүүдийг өргөтгөх/хураах
    \item Дээд/доод урсгалын шүүлтүүр
    \item Зангилааны хайлт болон тодруулалт
\end{itemize}

\paragraph{Нөлөөллийн шинжилгээ}
Өөрчлөлт хийхээс өмнө хэрэглэгчид өөрчлөлтүүдийг симуляцилж, нөлөөлөх бүх доод түвшний хөрөнгүүдийг дүрслэн харах боломжтой. Энэ нь өөрчлөлтүүдийг хэрэгжүүлэхээс өмнө бүрэн ойлгохыг баталгаажуулснаар үйлдвэрлэлийн ослуудыг бууруулдаг.

\paragraph{API-д түрүүлсэн зохиомж}
Бүх функц RESTful API-уудаар хандагдах боломжтой бөгөөд CI/CD дамжуулах хоолой, өгөгдлийн каталогууд, тусгай програмуудтай интеграцийг боломжтой болгоно.

\subsubsection{Байршуулалтын загвар}

Шийдэл нь янз бүрийн дэд бүтцийн шаардлагатай байгууллагуудад уян хатан байдлыг хангаж, байгууллага дээрх болон үүлэн байршуулалтыг (AWS, Azure, GCP) хоёуланг нь дэмждэг.

%-------------------------------------------------------------------------------
% 9. ШААРДЛАГЫН ЖАГСААЛТ
%-------------------------------------------------------------------------------
\subsection{Шаардлагын жагсаалт}
\label{subsec:dl_requirements}

Энэ хэсэгт Өгөгдлийн Гарал Үүслийн Системийн функциональ болон функциональ бус шаардлагууд, мөн таамаглалууд, хамаарлууд, тодорхойлогдсон эрсдэлүүдийг баримтжуулна.

\subsubsection{Шаардлагын тодорхойлолт}

\paragraph{Функциональ шаардлагууд}

\begin{table}[htbp]
\centering
\caption{Функциональ шаардлагууд}
\label{tab:functional_req}
\begin{tabular}{|c|p{10cm}|}
\hline
\textbf{ID} & \textbf{Шаардлага} \\
\hline
FR-01 & Систем нь гарал үүслийн харилцааг гаргаж авахын тулд SQL DDL болон DML мэдэгдлүүдийг задлан шинжлэх ёстой. \\
\hline
FR-02 & Систем нь баганын түвшний гарал үүслийн хяналтыг дэмжих ёстой. \\
\hline
FR-03 & Систем нь гарал үүслийн интерактив графын дүрслэл хангах ёстой. \\
\hline
FR-04 & Систем нь дээд болон доод урсгалын гарал үүслийн дамжуулалтыг дэмжих ёстой. \\
\hline
FR-05 & Систем нь үүрэгт суурилсан хандалтын хяналтыг (RBAC) хэрэгжүүлэх ёстой. \\
\hline
FR-06 & Систем нь PDF, PNG, JSON, CSV форматуудаар гарал үүслийг экспортлох ёстой. \\
\hline
FR-07 & Систем нь бүх метадатаар бүрэн текст хайлт хангах ёстой. \\
\hline
FR-08 & Систем нь бүх хэрэглэгчийн үйлдлийн аудит бүртгэл хөтлөх ёстой. \\
\hline
FR-09 & Систем нь автоматжуулсан метадата сканнерын хуваарь тохируулахыг дэмжих ёстой. \\
\hline
FR-10 & Систем нь гадаад интеграцид зориулсан RESTful API-уудыг хангах ёстой. \\
\hline
\end{tabular}
\end{table}

\paragraph{Функциональ бус шаардлагууд}

\begin{itemize}
    \item \textbf{Гүйцэтгэл:} Гарал үүслийн асуулгууд 10,000 хүртэлх зангилаа бүхий графуудад 2 секундын дотор үр дүнг буцаах ёстой.
    
    \item \textbf{Масштаблах чадвар:} Систем нь доройтолгүйгээр 1 сая хүртэлх өгөгдлийн хөрөнгийг дэмжих ёстой.
    
    \item \textbf{Бэлэн байдал:} Ажлын цагаар 99.5\% uptime зорилт.
    
    \item \textbf{Аюулгүй байдал:} Дамжуулалт дахь бүх өгөгдлийг TLS 1.2+ ашиглан шифрлэх ёстой.
    
    \item \textbf{Хэрэглэхэд хялбар байдал:} Шинэ хэрэглэгчид сургалтгүйгээр 5 минутын дотор гарал үүслийг удирдах чадвартай байх ёстой.
\end{itemize}

\subsubsection{Таамаглал ба хамаарлууд}

\paragraph{Таамаглалууд}
\begin{itemize}
    \item Эх өгөгдлийн сангууд системийн каталог болон мэдээллийн схемд унших хандалт хангадаг.
    \item SQL скриптүүд стандарт ANSI SQL синтакс эсвэл дэмжигдсэн диалектын хувилбаруудыг дагадаг.
    \item Хэрэглэгчид JavaScript идэвхжсэн орчин үеийн веб хөтчүүдтэй (Chrome, Firefox, Safari, Edge).
    \item Гарал үүслийн систем болон бүх хянагдаж буй өгөгдлийн эх үүсвэрүүдийн хооронд сүлжээний холболт байдаг.
\end{itemize}

\paragraph{Хамаарлууд}
\begin{itemize}
    \item Neo4j Community эсвэл Enterprise Edition (v4.4+)
    \item PostgreSQL (v13+)
    \item Node.js Runtime (v16+)
    \item Python (v3.9+) ANTLR runtime-тай
    \item Контейнержуулсан байршуулалтад Docker ба Kubernetes
\end{itemize}

\subsubsection{Эрсдэл}

\begin{table}[htbp]
\centering
\caption{Эрсдэлийн үнэлгээ}
\label{tab:risks}
\begin{tabular}{|p{4cm}|c|c|p{5cm}|}
\hline
\textbf{Эрсдэл} & \textbf{Магадлал} & \textbf{Нөлөө} & \textbf{Бууруулах арга} \\
\hline
Задлагч нарийн төвөгтэй SQL-д алдагдах & Дунд & Өндөр & Нөөц гарын авлагын гарал үүслийн оруулалт хэрэгжүүлэх; дүрмийн хамрах хүрээг өргөтгөх \\
\hline
Масштабд гүйцэтгэл буурах & Бага & Өндөр & Асуулгын оновчлол хэрэгжүүлэх; кэш давхарга нэмэх \\
\hline
Аюулгүй байдлын сул тал & Дунд & Өндөр & Тогтмол аюулгүй байдлын аудит; нэвтрэлтийн тест \\
\hline
Хэрэглэгчийн нэвтрүүлэлтийн эсэргүүцэл & Дунд & Дунд & UX дизайнд хөрөнгө оруулах; сургалт явуулах \\
\hline
Интеграцийн алдаа & Бага & Дунд & Иж бүрэн API тест; баримтжуулалт \\
\hline
\end{tabular}
\end{table}

%-------------------------------------------------------------------------------
% 10. UML ДИАГРАММУУД
%-------------------------------------------------------------------------------
\subsection{UML диаграммууд}
\label{subsec:dl_uml}

Энэ хэсэгт Өгөгдлийн Гарал Үүслийн Системийн бүтэц болон зан төлөвийг тодорхойлсон гол UML диаграммуудыг танилцуулна. Диаграмм бүрийг текстээр тайлбарлаж, хялбаршуулсан дүрслэлтэй хамт өгнө.

\subsubsection{Use Case диаграмм}

Use Case диаграмм нь үндсэн оролцогчид болон тэдний системтэй харилцааг тодорхойлно. Гурван үндсэн оролцогч Өгөгдлийн Гарал Үүслийн Системтэй харилцдаг: \textbf{Өгөгдлийн инженер} (эх үүсвэрүүдийг тохируулах, сканнер эхлүүлэх), \textbf{Өгөгдлийн аналист} (гарал үүслийг үзэх, хайлт хийх), \textbf{Администратор} (хэрэглэгчдийг удирдах, системийн тохиргоо).

Гол use case-үүд нь: Гарал үүсэл үзэх, Хөрөнгө хайх, Нөлөөллийн шинжилгээ харах, Өгөгдлийн эх үүсвэр тохируулах, Сканнер эхлүүлэх, Гарал үүсэл экспортлох, Хэрэглэгчдийг удирдах, Аудит бүртгэл харах. Өгөгдлийн инженер болон Аналист нь гарал үүсэл үзэх нийтлэг use case-үүдийг хуваалцдаг бол удирдлагын функцууд нь Администраторын үүрэгт хязгаарлагддаг.

\begin{verbatim}
+--------------------------------------------------+
|     Өгөгдлийн Гарал Үүслийн Систем                |
|                                                   |
|   (Гарал үүсэл үзэх)----+                        |
|   (Хөрөнгө хайх)--------+----> [Өг. Аналист]     |
|   (Нөлөөлөл харах)------+                        |
|                                                   |
|   (Эх үүсвэр тохируулах)+                        |
|   (Сканнер эхлүүлэх)----+----> [Өг. Инженер]     |
|   (Гарал үүсэл экспорт)-+                        |
|                                                   |
|   (Хэрэглэгч удирдах)---+                        |
|   (Аудит бүртгэл харах)-+----> [Администратор]   |
|   (Системийн тохиргоо)--+                        |
+--------------------------------------------------+
\end{verbatim}

\subsubsection{Class диаграмм}

Class диаграмм нь системийн үндсэн домэйн моделийг харуулна. Гол ангиуд нь \texttt{DataAsset} (бүх өгөгдлийн нэгжийн хийсвэр суурь), \texttt{Table}, \texttt{Column}, \texttt{ETLJob}, \texttt{Report} юм. \texttt{LineageEdge} анги нь хөрөнгүүдийн хоорондын харилцааг төлөөлж, хувиргалтын төрөл болон итгэлцлийн оноо зэрэг шинж чанаруудтай.

\texttt{User} анги нь хандалтын хяналтад \texttt{Role}-той холбогддог. \texttt{DataSource} нь тохируулагдсан гадаад системүүдийг төлөөлдөг бол \texttt{ScanJob} нь метадата гаргаж авалтын ажиллагааг хянадаг.

\begin{verbatim}
+----------------+       +------------------+
|   DataAsset    |<------|   LineageEdge    |
+----------------+       +------------------+
| -id: String    |       | -sourceId: String|
| -name: String  |       | -targetId: String|
| -type: String  |       | -transform: String|
| -metadata: Map |       +------------------+
+----------------+
       ^
       |
+------+-------+--------+
|              |        |
+--------+ +--------+ +--------+
| Table  | | ETLJob | | Report |
+--------+ +--------+ +--------+

+--------+     +--------+
|  User  |---->|  Role  |
+--------+     +--------+
\end{verbatim}

\subsubsection{Activity диаграмм}

Activity диаграмм нь гарал үүслийн гаргаж авалтын ажлын урсгалыг дүрсэлнэ. Үйл явц нь хэрэглэгч эсвэл хуваарилагч сканнер эхлүүлэх үед эхэлнэ. Систем нь тохируулагдсан өгөгдлийн эх үүсвэрт холбогдож, DDL/DML мэдэгдлүүдийг гаргаж авч, задлагчийн хөдөлгүүрт дамжуулна.

Задлагч нь оролтыг токенжуулж, AST (Abstract Syntax Tree) байгуулж, эх үүсвэр болон зорилтын нэгжүүдийг тодорхойлж, баганын зураглалыг гаргаж авна. Задлан шинжилсэн харилцаануудыг одоо байгаа метадататай харьцуулан баталгаажуулж, дараа нь граф өгөгдлийн санд хадгална. Эцэст нь шинэ гарал үүслийг тусгахын тулд UI кэшийг хүчингүй болгоно.

\begin{verbatim}
[Эхлэл]
   |
   v
(Сканнер эхлүүлэх)
   |
   v
(Өгөгдлийн эх үүсвэрт холбогдох)
   |
   v
(SQL/DDL гаргаж авах)
   |
   v
(Мэдэгдлүүдийг задлах)-----> <Шийдвэр: Задлалт амжилттай?>
   |                               |
   | Тийм                          | Үгүй
   v                               v
(Харилцаануудыг гаргах)       (Алдаа бүртгэх)
   |                               |
   v                               |
(Гарал үүслийг баталгаажуулах)     |
   |                               |
   v                               |
(Граф ӨС-д хадгалах)              |
   |                               |
   v                               |
(Кэш хүчингүй болгох) <-----------+
   |
   v
[Төгсгөл]
\end{verbatim}

\subsubsection{Sequence диаграмм}

Sequence диаграмм нь хэрэглэгч тодорхой хүснэгтийн гарал үүслийг хүсэх үеийн харилцааны урсгалыг харуулна. Хэрэглэгч UI-ээр хүсэлтийг эхлүүлж, UI нь API Gateway-г дууддаг. Gateway нь Auth Service-ээр хүсэлтийг баталгаажуулж, дараа нь асуулгыг Lineage Service руу дамжуулна.

Lineage Service нь Cypher асуулгууд ашиглан Neo4j-ээс дээд болон доод урсгалын харилцаануудыг асууна. Үр дүнг зангилаа болон ирмэгүүдийг агуулсан хариултын объект болгон угсарна. Хариулт нь Gateway-ээр UI руу буцаж очиж, интерактив графыг рендэрлэнэ.

\begin{verbatim}
Хэрэглэгч    UI        Gateway    AuthService   LineageService   Neo4j
 |           |            |            |              |            |
 |--Хүсэлт-->|            |            |              |            |
 |           |--API Call->|            |              |            |
 |           |            |--Баталгаа->|              |            |
 |           |            |<--Token OK-|              |            |
 |           |            |--Гарал үүсэл авах-------->|            |
 |           |            |            |              |--Асуулга-->|
 |           |            |            |              |<--Зангилаа-|
 |           |            |<--Хариулт-----------------|            |
 |           |<--JSON-----|            |              |            |
 |<--Рендэр--|            |            |              |            |
 |           |            |            |              |            |
\end{verbatim}

%-------------------------------------------------------------------------------
% 11. ХЭРЭГЖҮҮЛЭЛТ
%-------------------------------------------------------------------------------
\subsection{Хэрэгжүүлэлт}
\label{subsec:dl_implementations}

Энэ хэсэгт Өгөгдлийн Гарал Үүслийн Системийн гурван үндсэн модулийн хэрэгжүүлэлтийг дэлгэрэнгүй тайлбарлана: Хэрэглэгчийн модуль, Задлагч, Графикийн хэрэглэгчийн интерфэйс.

\subsubsection{Хэрэглэгчийн модуль}

Хэрэглэгчийн модуль нь Өгөгдлийн Гарал Үүслийн Системд иж бүрэн хэн мөн болохыг удирдах, баталгаажуулалт, зөвшөөрлийн чадамжуудыг хангадаг. Энэ нь байгууллагын шаталсан бүтэц болон үүрэгт суурилсан зөвшөөрлүүдийг дэмжихийн зэрэгцээ аюулгүй хандалтыг хангадаг.

\paragraph{Зорилго ба онцлогууд}
Модуль нь бүртгэлээс идэвхгүй болгох хүртэлх хэрэглэгчийн амьдралын мөчлөгийг удирдаж, профайл, итгэмжлэл, хандалтын зөвшөөрлүүдийг хадгалдаг. Гол функцууд нь:

\begin{itemize}
    \item Имэйл баталгаажуулалттай хэрэглэгчийн бүртгэл
    \item bcrypt хэшлэлт ашигласан аюулгүй нууц үг хадгалалт
    \item JWT токенуудаар сесс удирдлага
    \item Нууц үг сэргээх функц
    \item Профайл удирдлага (нэр, имэйл, тохиргоо)
    \item Аудит нийцлийн үйл ажиллагааны бүртгэл
\end{itemize}

\paragraph{Нэвтрэлт ба баталгаажуулалт}
Баталгаажуулалт нь байгууллагын нийцтэй байдлын OAuth 2.0 / OpenID Connect стандартуудыг дагадаг. Систем нь олон баталгаажуулалтын аргуудыг дэмждэг:

\begin{itemize}
    \item \textbf{Дотоод баталгаажуулалт:} Хэрэглэгчийн нэр/нууц үг, нэмэлт MFA (TOTP)-тай
    \item \textbf{SSO интеграци:} Байгууллагын хэн мөн болохын үйлчилгээ үзүүлэгчдэд зориулсан SAML 2.0
    \item \textbf{API баталгаажуулалт:} Программчлалын хандалтад зориулсан API түлхүүрүүд болон JWT bearer токенууд
\end{itemize}

JWT токенууд нь хэрэглэгчийн хэн мөн болох, үүрэг, дуусах хугацааны мэдээллийг агуулдаг. Хандалтын токенууд 1 цагийн дараа дуусдаг; шинэчлэх токенууд нь 7 хүртэлх хоногт саадгүй дахин баталгаажуулалтыг боломжтой болгодог.

\paragraph{Хэрэглэгчийн үүргүүд}
Систем нь шаталсан зөвшөөрөл бүхий гурван урьдчилан тодорхойлсон үүргийг хэрэгжүүлдэг:

\begin{itemize}
    \item \textbf{Үзэгч:} Гарал үүслийн графууд болон метадатад зөвхөн унших хандалт. Тохиргоог өөрчлөх эсвэл өгөгдөл экспортлох боломжгүй.
    
    \item \textbf{Засварлагч:} Үзэгчийн бүрэн зөвшөөрөл дээр гарын авлагын гарал үүслийн оруулга нэмэх, хөрөнгө тэмдэглэх, өгөгдөл экспортлох чадвартай.
    
    \item \textbf{Администратор:} Хэрэглэгчийн удирдлага, өгөгдлийн эх үүсвэрийн тохиргоо, системийн тохиргоо зэрэг системд бүрэн хандалт.
\end{itemize}

\paragraph{Гол функцууд}
Үндсэн API endpoint-ууд:
\begin{itemize}
    \item \texttt{POST /auth/login} --- Хэрэглэгчийн итгэмжлэлийг баталгаажуулах
    \item \texttt{POST /auth/refresh} --- Хандалтын токеныг шинэчлэх
    \item \texttt{GET /users/me} --- Одоогийн хэрэглэгчийн профайлыг авах
    \item \texttt{PUT /users/:id} --- Хэрэглэгчийн мэдээллийг шинэчлэх
    \item \texttt{GET /users/:id/activity} --- Хэрэглэгчийн үйл ажиллагааны бүртгэлийг авах
\end{itemize}

\subsubsection{Задлагч}

Задлагч модуль нь Өгөгдлийн Гарал Үүслийн Системийн аналитик хөдөлгүүр бөгөөд олон төрлийн өгөгдлийн эх үүсвэрүүдээс метадата болон гарал үүслийн харилцааг гаргаж авах үүрэгтэй. Энэ нь түүхий SQL, ETL скриптүүд, схемийн тодорхойлолтуудыг бүтэцлэгдсэн гарал үүслийн графууд болгон хувиргадаг.

\paragraph{Метадата гаргаж авах үйл явц}
Задлагч нь олон үе шаттай гаргаж авах дамжуулах хоолой ашигладаг:

\begin{enumerate}
    \item \textbf{Эх үүсвэрийн холболт:} Тохируулагдсан итгэмжлэл ашиглан зорилтот өгөгдлийн сан эсвэл файлын системд аюулгүй холболт үүсгэх.
    
    \item \textbf{Бүтээгдэхүүн цуглуулалт:} Системийн каталогуудаас (\texttt{information\_schema}), асуулгын бүртгэлүүд, файлын хадгалах сангуудаас DDL мэдэгдлүүдийг татаж авах.
    
    \item \textbf{Лексик шинжилгээ:} SQL диалект бүрт зориулсан ANTLR-ээр үүсгэгдсэн лексерүүд ашиглан оролтыг токенжуулах.
    
    \item \textbf{Синтакс задлан шинжлэл:} Асуулгын бүтцийг төлөөлсөн Abstract Syntax Trees (AST) байгуулах.
    
    \item \textbf{Семантик шинжилгээ:} Хүснэгт/баганын лавлагаануудыг шийдвэрлэх, alias-уудыг зохицуулах, хувиргалтуудыг тодорхойлох.
    
    \item \textbf{Гарал үүсэл гаргаж авалт:} Хувиргалтын логикоор дамжуулан эх баганыг зорилтын баганатай зураглах.
\end{enumerate}

\paragraph{Дэмжигдсэн форматууд}
Задлагч нь дараах оролтын форматуудыг дэмждэг:

\begin{itemize}
    \item \textbf{SQL диалектууд:} PostgreSQL, MySQL, Oracle PL/SQL, T-SQL (SQL Server), ANSI SQL
    \item \textbf{DDL мэдэгдлүүд:} CREATE TABLE, CREATE VIEW, ALTER TABLE
    \item \textbf{DML мэдэгдлүүд:} SELECT, INSERT, UPDATE, MERGE, CTAS
    \item \textbf{ETL тодорхойлолтууд:} Apache Airflow DAG-ууд (Python), dbt моделүүд (SQL + YAML)
    \item \textbf{Схем файлууд:} JSON Schema, Avro, Parquet метадата
\end{itemize}

\paragraph{Гаралтын формат}
Задлан шинжилсэн гарал үүсэл нь стандартчилагдсан схемд нийцсэн JSON баримт бичиг хэлбэрээр гарна:

\begin{lstlisting}[language=JavaScript, caption=Гарал үүслийн гаралтын схем]
{
  "nodes": [
    {
      "id": "source_db.schema.table",
      "type": "TABLE",
      "columns": ["col1", "col2"]
    }
  ],
  "edges": [
    {
      "source": "source_db.schema.table.col1",
      "target": "target_db.schema.table.col_a",
      "transformation": "DIRECT_COPY",
      "confidence": 1.0
    }
  ]
}
\end{lstlisting}

\subsubsection{Графикийн хэрэглэгчийн интерфэйс}

Графикийн хэрэглэгчийн интерфэйс (GUI) нь гарал үүслийн өгөгдөлтэй харилцах зөн совинтой, веб-д суурилсан frontend-ийг хангадаг. Энэ нь гарын товчлол болон дэвшилтэт шүүлтүүрээр дамжуулан хүчирхэг хэрэглэгчийн ажлын урсгалыг дэмжихийн зэрэгцээ харааны судлалыг онцолдог.

\paragraph{Фреймворк ба технологи}
GUI нь орчин үеийн веб технологиуд ашиглан бүтээгдсэн:

\begin{itemize}
    \item \textbf{React 18:} Төлөв удирдлагад hooks бүхий бүрэлдэхүүн хэсэгт суурилсан UI фреймворк
    \item \textbf{TypeScript:} Сайжруулсан засвар үйлчилгээнд статик тайпчлал
    \item \textbf{D3.js:} Гарал үүслийн графуудад зориулсан өгөгдөлд суурилсан дүрслэлийн сан
    \item \textbf{Material-UI:} Тогтвортой дизайны хэлийг хангадаг бүрэлдэхүүн хэсгийн сан
    \item \textbf{React Query:} Кэшлэлттэй серверийн төлөв удирдлага
\end{itemize}

\paragraph{Үндсэн дэлгэцүүд}
Интерфэйс нь хэд хэдэн гол дэлгэцээс бүрдэнэ:

\begin{itemize}
    \item \textbf{Хянах самбар:} Сүүлийн сканнерууд, гарал үүслийн статистик, хурдан хайлтын тойм.
    
    \item \textbf{Каталог хөтөч:} Эх үүсвэр, схем, төрлөөр өгөгдлийн хөрөнгүүдийн шаталсан навигаци.
    
    \item \textbf{Гарал үүслийн судлагч:} Интерактив графтай үндсэн дүрслэлийн дэлгэц.
    
    \item \textbf{Хайлтын үр дүн:} Хөрөнгийн төрөл, эх үүсвэр, шошгоор талын шүүлтүүр бүхий бүрэн текст хайлт.
    
    \item \textbf{Хөрөнгийн дэлгэрэнгүй:} Бие даасан хөрөнгийн метадата, гарал үүсэл, хэрэглээний статистикийн иж бүрэн харагдац.
    
    \item \textbf{Удирдлага:} Хэрэглэгчийн удирдлага, эх үүсвэрийн тохиргоо, системийн тохиргоо.
\end{itemize}

\paragraph{Гарал үүслийн графын дүрслэл}
GUI-ийн голд интерактив гарал үүслийн граф байна. Функцууд нь:

\begin{itemize}
    \item Органик зангилааны байршуулалтад физикийн симуляци бүхий хүчний чиглэлтэй зохион байгуулалт
    \item Томруулах (хулганы дугуй) болон гүйлгэх (чирэх) навигаци
    \item Холбогдсон зангилаануудыг аажмаар илрүүлэхийн тулд дарж өргөтгөх
    \item Хөрөнгийн төрлөөр өнгөний код (хүснэгтүүд = цэнхэр, view-үүд = ногоон, ETL = улбар шар)
    \item Хувиргалтын төрлүүдийг харуулсан ирмэгийн шошго
    \item Том графуудад чиглүүлэхэд зориулсан жижиг зураг
\end{itemize}

\paragraph{Интерактив функцууд}
Хэрэглэгчид графтай дараах байдлаар харилцаж болно:
\begin{itemize}
    \item \textbf{Төвлөрлийн горим:} Сонгосон зангилаанаас дээд эсвэл доод урсгалын замуудыг тодруулах
    \item \textbf{Шүүлтүүрийн самбар:} Төрөл, эх үүсвэр, хайлтын шалгуураар зангилаануудыг харуулах/нуух
    \item \textbf{Дэлгэрэнгүй хажуугийн самбар:} Графаас гарахгүйгээр сонгосон зангилааны метадатаг харах
    \item \textbf{Тэмдэглэгээ:} Хамтын ажиллагаанд зориулж зангилаанд тайлбар болон шошго нэмэх
\end{itemize}

\paragraph{Экспортын сонголтууд}
Гарал үүслийг баримтжуулалт болон хуваалцахад экспортлох боломжтой:
\begin{itemize}
    \item \textbf{PNG/SVG:} Одоогийн графын харагдацын өндөр нарийвчлалтай зураг экспорт
    \item \textbf{PDF:} Гарал үүслийн граф болон метадатаг багтаасан форматлагдсан тайлан
    \item \textbf{JSON:} Программчлалын хэрэглээнд зориулсан машин уншигдах гарал үүсэл
    \item \textbf{CSV:} Зангилаа болон ирмэгийн жагсаалтын хүснэгт экспорт
\end{itemize}

%-------------------------------------------------------------------------------
% 12. ДҮГНЭЛТ
%-------------------------------------------------------------------------------
\subsection{Дүгнэлт}
\label{subsec:dl_conclusion}

Энэ тайланд танилцуулсан Өгөгдлийн Гарал Үүслийн Систем нь орчин үеийн өгөгдлийн удирдлагын чухал хэрэгцээг хангадаг: нарийн төвөгтэй байгууллагын орчинд өгөгдлийн урсгалыг автоматжуулсан хяналт. Ухаалаг задлан шинжлэл, граф-д суурилсан хадгалалт, зөн совинтой дүрслэлийг хослуулснаар систем нь тодорхойгүй өгөгдлийн экосистемүүдийг ил тод, навигацилах боломжтой орчин болгон хувиргадаг.

\subsubsection{Амжилтуудын товч}

Энэ төсөл амжилттайгаар дараах зүйлсийг хүргүүлсэн:

\begin{itemize}
    \item Олон өгөгдлийн эх үүсвэрүүдээс автоматжуулсан гарал үүслийн гаргаж авалтыг дэмждэг иж бүрэн архитектур.
    
    \item Гол өгөгдлийн сангийн платформуудаар баганын түвшний нарийвчлалтай SQL-ийг шинжлэх чадвартай боловсронгуй задлагчийн хөдөлгүүр.
    
    \item Гарал үүслийг техникийн болон техникийн бус хэрэглэгчдэд хоёуланд нь хүртээмжтэй болгодог интерактив графын дүрслэл бүхий зөн совинтой веб интерфэйс.
    
    \item Байгууллагын байршуулалтад тохирсон аюулгүй, үүрэгт суурилсан хандалтын хяналтын систем.
    
    \item Одоо байгаа өгөгдлийн дэд бүтэц болон CI/CD дамжуулах хоолойтой интеграцийг боломжтой болгодог RESTful API-ууд.
\end{itemize}

\subsubsection{Бизнесийн үнэ цэнэ}

Энэ системийг хэрэгжүүлж буй байгууллагууд дараах зүйлсийг хүлээж болно:
\begin{itemize}
    \item Гарын авлагын гарал үүслийн баримтжуулалтад зарцуулах цагийг 60-70\%-иар бууруулах
    \item Аудит хийгдэх боломжтой гарал үүслийн бүртгэлээр дамжуулан зохицуулалтын нийцлийг сайжруулах
    \item Идэвхтэй нөлөөллийн шинжилгээгээр дамжуулан үйлдвэрлэлийн ослуудыг бууруулах
    \item Аналитик санаачилгуудыг хурдасгах сайжруулсан өгөгдөл олох чадвар
\end{itemize}

\subsubsection{Ирээдүйн сайжруулалтууд}

Дараах сайжруулалтуудыг ирээдүйн хувилбаруудад төлөвлөж байна:

\begin{enumerate}
    \item \textbf{Бодит цагийн гарал үүсэл:} Kafka, Spark Streaming болон ижил төстэй платформуудаас гарал үүсэл барихад зориулсан урсгал боловсруулалтын интеграци.
    
    \item \textbf{Машин сургалтын сайжруулалт:} Баримтжуулалт муутай хуучин системүүдэд зориулсан AI-д туслагдсан гарал үүслийн дүгнэлт.
    
    \item \textbf{Өгөгдлийн чанарын интеграци:} Асуудал автоматаар тархах сэрэмжлүүлэг бүхий гарал үүсэлд мэдрэмтгий өгөгдлийн чанарын хяналт.
    
    \item \textbf{BI хэрэгслийн холбогчид:} Тайлангийн түвшний гарал үүсэлд зориулсан Tableau, Power BI, Looker-тэй уугуул интеграци.
    
    \item \textbf{Хамтын ажиллагааны функцууд:} Гарал үүслийн өөрчлөлтөд зориулсан тайлбар, мэдэгдэл, ажлын урсгалын батлал.
    
    \item \textbf{Гар утасны програм:} Явж байхдаа гарал үүсэлд хандахад зориулсан iOS болон Android програмууд.
\end{enumerate}

Өгөгдлийн Гарал Үүслийн Систем нь иж бүрэн өгөгдлийн засаглалын суурийг бий болгож, байгууллагуудыг өөрсдийн өгөгдлийн хөрөнгийг итгэлтэй, тодорхой байдлаар удирдах боломжийг олгоно.

%-------------------------------------------------------------------------------
% 13. НОМ ЗҮЙ
%-------------------------------------------------------------------------------
\subsection{Ном зүй}
\label{subsec:dl_references}

\begin{enumerate}
    \item Thompson, R., \& Garcia, M. (2019). Байгууллагын өгөгдлийн засаглалд зориулсан автоматжуулсан метадата удирдлагын фреймворкууд. \textit{Journal of Data Management}, 15(3), 234--251.
    
    \item Chen, L., Wang, H., \& Zhang, Y. (2020). Нарийн төвөгтэй өгөгдлийн гарал үүслийн системүүдэд зориулсан граф-д суурилсан дүрслэлийн техникүүд. \textit{IEEE Transactions on Visualization and Computer Graphics}, 26(8), 2891--2904.
    
    \item Williams, J., \& Brown, S. (2018). Байгууллагын өгөгдлийн гарал үүслийн интеграцийн сорилтууд: Олон платформын судалгаа. \textit{International Journal of Information Management}, 42, 78--92.
    
    \item Kumar, A., \& Patel, R. (2021). Автомат SQL гарал үүсэл гаргаж авахад машин сургалтын аргууд. \textit{Proceedings of the ACM SIGMOD Conference}, 1156--1168.
    
    \item Anderson, K. (2020). GDPR болон CCPA-ийн эрин дэх өгөгдлийн гарал үүслийн зохицуулалтын үр нөлөө. \textit{Journal of Data Protection \& Privacy}, 4(2), 112--128.
    
    \item Liu, X., \& Yamamoto, T. (2022). Урсгал өгөгдлийн архитектурт үйл явдалд суурилсан гарал үүсэл барих. \textit{Proceedings of the IEEE International Conference on Big Data}, 892--901.
    
    \item Martinez, C. (2019). \textit{Өгөгдлийн засаглал: Орчин үеийн байгууллагын зарчим ба практик}. O'Reilly Media.
\end{enumerate}

\vfill
